\documentclass[11pt]{article}
\usepackage{ctex}
\usepackage{geometry}
\usepackage{indentfirst}
\usepackage{titlesec}
\usepackage{enumitem}
\usepackage{fancyhdr}
\usepackage{amsmath}
\usepackage{amsfonts}
\usepackage{amssymb}

% 设置页面边距
\geometry{a4paper, left=2.5cm, right=2.5cm, top=2.5cm, bottom=2.5cm}

% 设置页眉页脚
\pagestyle{fancy}
\fancyhf{}
\fancyhead[L]{南京大学}
\fancyhead[C]{自然语言处理大项目}
\fancyhead[R]{\today}
\fancyfoot[C]{\thepage}
\renewcommand{\headrulewidth}{0.4pt}

% 设置标题格式
\titleformat{\section}{\large\bfseries}{\thesection}{1em}{}
\titleformat{\subsection}{\normalsize\bfseries}{\thesubsection}{1em}{}

% 设置段落间距
\setlength{\parindent}{2em} % 首行缩进
\setlength{\parskip}{0.5em}  % 段落间距
\setlength{\headheight}{13.6pt}

\title{\textbf{NLP大项目Proposal：长期记忆对话Agent}}
\author{宋昌致、毛泽毅、陈柯男}
\date{2026年5月17日}

\begin{document}

\maketitle

\section{选题}

本项目旨在实现一个支持长期记忆的对话Agent，该Agent能够从历史多轮对话中自动提取并管理记忆，在回答新问题时检索并利用相关记忆，并在长期对话评测集上测量效果。

具体来说，我们将构建一个系统，该系统包含以下核心模块：

\begin{enumerate}
\item Memory Store：记忆存储与索引（使用Chroma）
\item Memory Writer：从对话中提取记忆（基于LLM的prompt）
\item Memory Retriever：给定query检索相关记忆
\item Memory Updater：实现冲突机制
\item Agent Controller：主流程编排：写入 → 检索 → 生成
\end{enumerate}

我们的方案将参考以下四篇核心文献的思路：

\begin{enumerate}
\item Generative Agents: Interactive Simulacra of Human Behavior (Park et al., UIST 2023) - Recency × Importance × Relevance 三因子检索、Reflection
\item MemoryBank: Enhancing Large Language Models with Long-Term Memory (Zhong et al., AAAI 2024) - 基于 Ebbinghaus 遗忘曲线的记忆动力学
\item Mem0: Building Production-Ready AI Agents with Scalable Long-Term Memory (Chhikara et al., ECAI 2025) - 记忆提取 → 冲突检测 → 更新的工程化流程
\item Evaluating Very Long-Term Conversational Memory of LLM Agents (Maharana et al., ACL 2024) - 长期对话记忆评测基准
\end{enumerate}

\section{组队情况}

\begin{itemize}
\item 队长：宋昌致 241880436
\item 队员：毛泽毅 241880422、陈柯男 241880092
\end{itemize}

\section{研究计划}

\subsection{项目总体目标}

实现一个支持长期记忆的对话Agent，具备以下能力：

\begin{enumerate}
\item 从多会话对话中提取并结构化存储记忆
\item 回答新问题时检索相关记忆辅助生成
\item 实现基本的记忆管理机制——选择冲突机制作为本团队核心目标
\item 在评测集上跑通并得到可量化的结果
\end{enumerate}

\subsection{技术路线}

我们将采用以下技术栈：

\begin{enumerate}
\item 主LLM：Qwen2.5-3B-Instruct-AWQ
\item Embedding：BAAI/bge-m3
\item 部署方式：使用vLLM起OpenAI兼容服务
\item 评测：使用云端API（DashScope qwen-plus）进行LLM-as-Judge评测
\end{enumerate}

\subsection{探索方向}

根据课程要求，我们选择探索方向A（检索策略）和E（反思机制），其中以E（反思机制）为主要探索方向，如有时间再探索A（检索策略）。

\begin{enumerate}
\item \textbf{主要探索方向E：反思机制（Reflection）}
\begin{itemize}
\item 借鉴Generative Agents中的Reflection机制
\item 实现周期性reflection生成高阶记忆
\item 通过反思机制整合低层次观察形成高层次洞察
\end{itemize}

\item \textbf{次要探索方向A：检索策略（如有时间）}
\begin{itemize}
\item 对比稠密检索vs.三因子打分(Recency × Importance × Relevance)
\item 基于Generative Agents的检索策略实现
\end{itemize}
\end{enumerate}

\subsection{重要时间节点}

\begin{itemize}
\item \textbf{2026年5月17日：}Proposal 提交选题、组队情况、研究计划
\item \textbf{2026年6月3日：}Milestone 提交进展报告（进度、初步成果、下一步计划）
\item \textbf{2026年6月17日：}Final 答辩PPT、代码、论文
\end{itemize}

\subsection{详细实施计划}

\subsubsection{第1周（5月17日前）：Proposal完成与系统架构搭建}

\begin{enumerate}
\item 完成Proposal文档（截至5月17日）
\item 搭建可端到端运行的基线Agent框架
\item 实现Memory Store模块（使用FAISS或Chroma）
\item 实现Memory Writer模块（从对话中提取记忆）
\end{enumerate}

\subsubsection{第2周（5月18日-5月24日）：核心功能继续完善}

\begin{enumerate}
\item 实现Memory Retriever模块（检索相关记忆）
\item 实现Memory Updater模块（选择更新/冲突/遗忘中的一种）
\item 实现Agent Controller主流程编排
\end{enumerate}

\subsubsection{第3-4周（5月25日-6月1日）：进展报告准备与初步成果实现}

\begin{enumerate}
\item 确保所有记忆操作可追踪（记录检索到的记忆内容、完整prompt和模型输出）
\item 开始实现主要探索方向E：反思机制（Reflection）
\item 准备进展报告（截至6月3日前提交）
\begin{itemize}
\item 项目进度总结
\item 初步成果展示
\item 下一步详细计划
\end{itemize}
\end{enumerate}

\subsubsection{第5-6周（6月2日-6月15日）：主要探索方向深入实现}

\begin{enumerate}
\item 完善Reflection机制设计与实现
\item 实现周期性高阶记忆生成
\item 开发记忆整合算法，将低层次观察转化为高层次洞察
\item 进行初步的消融实验，验证Reflection机制的有效性
\item 如有时间，开始次要探索方向A：检索策略实现
\begin{itemize}
\item 实现多种检索策略（稠密检索vs.三因子打分）
\item 对比不同检索策略的效果
\end{itemize}
\end{enumerate}

\subsubsection{第7周（6月16日）：答辩准备}

\begin{enumerate}
\item 完成所有功能开发
\item 进行全面测试和性能评估
\item 准备Final答辩PPT
\item 整理项目代码和论文
\end{enumerate}

\subsubsection{2026年6月17日：Final答辩}

\begin{enumerate}
\item 提交答辩PPT
\item 提交完整项目代码
\item 提交项目论文
\item 进行项目答辩展示
\end{enumerate}

\subsection{预期成果}

\begin{enumerate}
\item 一个可运行的长期记忆对话Agent系统，具备反思机制
\item 完整的实验评估结果，包括QA准确率、LLM调用次数、响应时间等指标
\item 消融实验结果，证明反思机制的有效性
\item 详细的项目报告，包含系统架构图、实验设置、结果分析和反思
\item 6月3日前的进展报告
\item 6月17日的答辩PPT
\end{enumerate}

\subsection{风险评估与应对措施}

\begin{enumerate}
\item \textbf{风险1：}时间紧迫，难以在截止日期前完成所有功能 \\
      \textbf{应对：}严格按照时间节点推进，优先完成核心功能，确保在各截止日期前交付相应成果

\item \textbf{风险2：}反思机制实现复杂度高 \\
      \textbf{应对：}先实现基础版本，逐步迭代优化，确保在6月3日前有初步成果展示

\item \textbf{风险3：}评测数据集处理复杂 \\
      \textbf{应对：}先用小样本（5-10条）调通整条流水线，再扩展到全集评测

\item \textbf{风险4：}次要探索方向（检索策略）时间不足 \\
      \textbf{应对：}优先保证主要探索方向（反思机制）的完成质量，次要方向根据时间灵活调整
\end{enumerate}

\end{document}